\section{Introduction}

Large-scale convex optimization deals with problems whose dimension
\(n\) is huge, yet we require iterative algorithms whose number of
iterations is \(\ll n\). Convexity is one of the very few structural
properties that render such problems “solvable” both theoretically
and practically.

\subsection{Mathematical Optimization Problem}
\begin{equation}
  \begin{aligned}
    \operatorname{minimize}\quad & f(x)\\
    \text{s.t.}\quad & g_i(x)\le 0,\quad i=1,\dots,n_g,\\
    & h_i(x)=0,\quad i=1,\dots,n_h.
  \end{aligned}
  \label{eq:opt}
\end{equation}
- \(x\in\mathbb{R}^n\) – decision variable (algorithms also work for
\(n\to\infty\));
- \(f\) – objective;
- \(\mathcal{C}=\{\xi\in\mathbb{R}^n:g(\xi)\le 0,\,h(\xi)=0\}\) – feasible set.

\subsection{Existence of a Minimizer}
A minimizer exists when \(f\) is lower-semicontinuous,
\(f(x)\to+\infty\) as \(|x|\to\infty\), and \(\mathcal{C}\) is
closed. In this case we write
\[
  \min_{x\in\mathcal{C}}f(x)\qquad\text{and}\qquad
  x^*\in\arg\min_{x\in\mathcal{C}}f(x).
\]

\subsection{First-Order Iterative Algorithms}
All methods considered in the course are first-order:
\[
  x_{k+1}\;\leftarrow\;\text{function of}\;\bigl\{(f(x_j),\nabla
  f(x_j))\bigr\}_{j=0}^k.
\]
The analytic complexity (number of oracle calls needed to reach
accuracy \(\varepsilon\)) is the central quantity; arithmetic
complexity (total floating-point operations) is secondary.

\subsection{Curse of Dimensionality}
For the class of \(L^\infty\)-Lipschitz problems on the unit
hypercube, any algorithm may require
\[
  N\ge\Bigl(\Bigl\lfloor\frac{L}{2\varepsilon}\Bigr\rfloor\Bigr)^n-1
\]
oracle calls. Hence, without convexity or other structure, even
modest accuracy is impossible for large \(n\).

\subsection{Convexity}
Problem~\eqref{eq:opt} is convex when \(f\) and all \(g_i\) are
convex and all \(h_i\) are affine. In the convex case every local
minimizer is global, which is the key property that makes large-scale
problems tractable.

